\documentclass[../main.tex]{subfiles}
\ifSubfilesClassLoaded{\setcounter{chapter}{15}}{}
\begin{document}

\chapter{Proyek Mini End-to-End (Draft)}

\begin{subcpmk}
  \item \textbf{Sub-CPMK-P10:} Menyusun draft alur data mentah $\rightarrow$ prapemrosesan $\rightarrow$ model $\rightarrow$ evaluasi $\rightarrow$ insight (RPS Mg.\ 15).
\end{subcpmk}

\SesiRPS{15}{P10}{$\sim$170 menit}{CPMK-P2, CPMK-P3, CPMK-P4}

\section{Lingkup proyek}

\begin{konsep}
Proyek mini membuktikan integrasi capaian P2--P4. Topik disetujui dosen: masalah jelas, data legal, minimal satu model, metrik selaras jenis target, kesimpulan non-teknis.
\end{konsep}

\section{Checklist draft (isi di notebook)}

\begin{enumerate}[leftmargin=*]
  \item Rumusan masalah dan pemangku kepentingan (CRISP-DM fase 1).
  \item Sumber data + lisensi + ukuran + variabel kunci.
  \item Prapemrosesan terdokumentasi (missing, encoding, split).
  \item Model dan baseline; metrik pada hold-out atau CV.
  \item Satu visual utama dan paragraf rekomendasi.
\end{enumerate}

\section{Contoh kerangka kode orkestrasi}

\begin{lstlisting}[caption=Kerangka import dan split proyek]
import pandas as pd
from sklearn.model_selection import train_test_split

# df = pd.read_csv("data_proyek.csv")  # ganti
# X, y = ...  # definisikan target
# X_train, X_test, y_train, y_test = train_test_split(...)
\end{lstlisting}

\begin{lstlisting}[caption=Pipeline disarankan]
from sklearn.pipeline import Pipeline
from sklearn.preprocessing import StandardScaler
from sklearn.linear_model import LogisticRegression

model = Pipeline([
    ("sc", StandardScaler()),
    ("clf", LogisticRegression(max_iter=1000)),
])
# model.fit(X_train, y_train)
\end{lstlisting}

\begin{lstlisting}[caption=Metrik dan simpan figur]
from sklearn.metrics import classification_report
# print(classification_report(y_test, model.predict(X_test)))
# plt.savefig("proyek_metrik.png")
\end{lstlisting}

\begin{lstlisting}[caption=Sumber terbuka dataset]
# Lihat https://data.go.id/ dan https://www.kaggle.com/
\end{lstlisting}

\section{Referensi}

\begin{itemize}[leftmargin=*]
  \item data.go.id. \url{https://data.go.id/}
  \item INRIA sklearn MOOC (proyek). \url{https://inria.github.io/scikit-learn-mooc/}
\end{itemize}

\begin{asesmen}
\textbf{Proyek parsial:} unggah \texttt{Proyek\_draft\_NIM.ipynb} + outline laporan jika diminta.
\end{asesmen}

\begin{rangkuman}
Minggu 15 mengonsolidasikan pipeline penuh dalam satu artefak utama.

\textbf{Kata kunci:} end-to-end, CRISP-DM, dokumentasi
\end{rangkuman}

\end{document}
